\subsection{Small extreme corners and relative triangle obstructions}

\begin{lemma}[Proved: relative small-band triangles after exact endpoint orbits]
\label{CP2:A-relative-small-band-endpoint-orbits}
Let $C$ have full linear support $0<1<2$, with corners $A,M,N$
and nonempty strict hom-sets $X=C(0,1)$, $Y=C(1,2)$ and
$Z=C(0,2)$. Suppose the trivial left $\mathbb QA$-module and
the trivial right $\mathbb QN$-module are projective. Put
\[
 r=|X/A|,\qquad s=|N\backslash Y|.
\]
If $BC$ has finite CW homotopy type and either
$\min(r,s)\le2$ or $r=s=3$, every rational product of two
degree-two classes on $BC$ is zero. No finite CW hypothesis
is imposed on the endpoint-orbit quotient.
\end{lemma}
\begin{proof}
The two actual extreme-orbit functors of
Lemma~\ref{CP2:A-high-orbits} give a category $T$ with trivial
endpoint corners, unchanged middle corner $M$, adjacent hom-sets
$V=X/A$, $W=N\backslash Y$, and long hom-set
$N\backslash Z/A$. The map $BC\to BT$ is a rational homology
equivalence and preserves products. The actual long-arrow
collapse $T\to\overline T$ has surjective pullback on $H^2$,
by Lemma~\ref{CP2:A-triangle-long-collapse}. This conclusion
about the map does not require a finite CW hypothesis on $BT$.

The compression in Lemma~\ref{CP2:A-singleton-triangle-monoid}
uses the literal monoid
\[
 P=M\sqcup(V\times W),
 \qquad (v,w)(v',w')=(v,w').
\]
More precisely, it supplies an actual functor from
$\overline T$ to the full two-object subcategory $D$ of the
Karoubi envelope of $P$, on the global identity and one band
idempotent. This functor and $P\hookrightarrow D$ are homotopy
equivalences. The middle-corner map is the literal inclusion
$M\hookrightarrow P$. Consequently every degree-two class on
$BC$ is pulled back from $BP$, through these specified maps.

If one band dimension is one, $H^2(BP;\mathbb Q)=0$. If one
dimension is two, every balanced matrix has rank at most one.
Lemma~\ref{A:rank-one-square} and rational polarization kill
all products on $BP$, without a boundedness hypothesis. These
cases therefore pull back to the desired conclusion.

Suppose $r=s=3$. Use the complete thirteen-orbit relation-hull
classification in Theorem~\ref{CP2:A-three-by-three-all-products}.
The actual action homomorphism $P\to H(R)$ fixes the minimal
band and induces an isomorphism on rational $H^2$. In every
hull except its character-pruned orbit $11$, that proof kills
all products on the hull itself, without any finite CW
assumption on the source or target. Those cases again pull back.

Here is the needed source argument for orbit $11$. Its sole
integral balanced class $u$ has nonzero cyclic restriction at
each forbidden actor. All band actors are allowed. Thus if a
forbidden actor belonged to the image of $P$, one of its
preimages would belong to the unchanged middle monoid $M$.
The eventual cyclic group of that preimage surjects onto the
eventual cyclic group of the forbidden actor. It is a local
group at an idempotent of $M$, hence is represented by an actual
object of the Karoubi envelope of the original category $C$.

Extend the preceding functors to Karoubi envelopes. Because
the middle map was the inclusion of $M$, the restriction of
the pullback of $u$ to this local group is exactly inflation
of its checked nonzero cyclic character. That inflation is
nonzero. On the other hand, the Karoubi envelope has classifying
space homotopy equivalent to $BC$, which has finite CW type.
The line-bundle nilpotence argument of
Proposition~\ref{A:three-row-classification}, equivalently the
source argument in Lemma~\ref{CP2:A-relative-three-by-three},
forbids precisely this nonzero local-group restriction.
It is applied to the original $BC$, not to $BT$ or $BP$.

The image of $P$ therefore avoids all forbidden actors.
Adjoining the whole band and the identity still avoids them,
because a product involving a band element lies in the band.
This image is contained in one of the nine checked maximal
allowed submonoids, each with zero rational $H^4$.
The homomorphism from $P$ factors through that target; its
extension to the two-object Karoubi subcategory gives the
corresponding actual factorization of the map from $C$.
All degree-two classes in question lift from $H(R)$, so their
restrictions to this smaller target have zero products and
their pullbacks to $BC$ do also. This proves the lemma.
\end{proof}

\begin{corollary}[Proved: a nontrivial order-two endpoint on a short stalk]
\label{CP2:A-linear19-two-endpoint-short-stalk}
A full-linear category of at most nineteen arrows cannot have
type $\CP^2$ if its endpoint orders satisfy $a\le2$, $c=2$,
its long hom-set is a singleton, and
\[
 y\le3\qquad\hbox{or}\qquad x\le3,\ y\le4.
\]
The opposite conditions give the same conclusion. Also, at
any arrow count, $a,c\le2$ together with either
$\min(x,y)\le2$ or $x,y\le3$ gives zero degree-two products
whenever $BC$ has finite CW type.
\end{corollary}
\begin{proof}
Both monoids of order two have semisimple rational algebras.
If the last corner acts nontrivially on $Y$, its orbit set has
at most $|Y|-1$ elements. In the first displayed case this is
at most two; in the second it is at most three and
$|X/A|\le3$. The preceding lemma applies. If that action is
trivial, its action on the singleton $Z$ is trivial as well.
The entire extreme action is then trivial, so
Theorem~\ref{CP2:A-trivial-extreme-corner-retract} gives an
actual retract of type $\CP^2$ with one fewer arrow.
The completed exclusion through eighteen arrows,
Theorem~\ref{CP2:A-global-lower-nineteen}, gives a contradiction.
The last assertion follows directly from the preceding lemma,
since exact endpoint orbits cannot increase either stalk size.
Opposition proves the dual statements.
\end{proof}

\begin{lemma}[Proved: a three-element sink on a ternary stalk]
\label{CP2:A-linear19-three-sink-ternary-stalk}
Let $C$ have at most nineteen arrows and full linear support.
If its initial corner is trivial, its last corner has order
at most three, $x,y\le3$, and $z=1$, then $BC$ cannot have
type $\CP^2$.
\end{lemma}
\begin{proof}
If the trivial right module of the last corner $N$ is
projective, Lemma~\ref{CP2:A-relative-small-band-endpoint-orbits}
applies, since both orbit stalks have at most three elements.
The sole remaining monoid through order three is
\[
 R_3=\{1,e_0,e_1\},\qquad e_i e_j=e_j.
\]
Indeed, at a minimal idempotent $e$, the projective-side
criterion is $eN=eNe$, with rational averaging over the local
group. Failure requires at least two minimal columns. The
minimal ideal is proper, so the size bound forces exactly one
row, two columns, trivial local group, and just the global
identity outside it. This is exactly $R_3$.

Write $P_i:Y\to Y$ for the left actions of $e_i$. They are
idempotent and $P_iP_j=P_j$, so they have the same nonempty
image $U$ and fix $U$ pointwise. Every commuting right-$M$
actor preserves $U$. If $|U|=1$, its point is fixed by both
$N$ and $M$. The joint fixed-middle-arrow lemma,
Lemma~\ref{CP2:A-linear-joint-fixed-middle-arrow}, kills all
degree-two products. If $U=Y$, the entire sink action on $Y$
and the singleton $Z$ is trivial. Its actual corner retract
strictly reduces the arrow count and contradicts the
eighteen-arrow exclusion.

Only $|Y|=3$, $|U|=2$ remains. If $P_0=P_1=P$, use the
idempotent cut map that is $P$ on $Y$ and the identity on $Z$.
It commutes with $N$ and $M$. The sole compatibility involving
$X$ holds because every two-cross-arrow composite lies in
the singleton $Z$. Thus this is an actual proper cut retract
of the full category. The cut-retract theorem transfers
type $\CP^2$ to its image, which has one fewer arrow, again
contradicting the eighteen-arrow exclusion.

If $P_0\ne P_1$, write $Y=\{u_0,u_1,w\}$ with
$P_iw=u_i$. Its left-$R_3$ action category has corner $R_3$
at each $u_i$, identity corner at $w$, and just the two strict
arrows $w\to u_i$ labelled $e_i$. The actual functor that
contracts the two $R_3$ corners to identities maps it to the
barycentric interval. Its comma categories are exactly those
checked in Lemma~\ref{CP2:A-directed-graph-right-zero-realization},
so this action category is contractible.

Now include the full last corner into $C$. The comma category
over the middle object is the displayed action category on
$Y$. The comma category over the initial object is the
action category on the singleton $Z$, namely $R_3$, also
contractible. At the last object the comma category has
initial object $1_N$. Quillen's Theorem~A therefore makes
$BN\to BC$ a homotopy equivalence. Since $BR_3$ is
contractible, this case cannot have projective-plane type.
The four cases exhaust every action on at most three states.
\end{proof}

\begin{corollary}[Proved: five further large middle-corner profiles]
\label{CP2:A-linear19-five-small-extreme-large-profiles}
The following nineteen-arrow profiles are excluded, together
with their opposites:
\[
\begin{gathered}
 (1,8,2;3,3,2),\quad (1,8,2;3,4,1),\quad
 (1,8,2;4,3,1),\\
 (1,8,3;3,3,1),\quad (2,8,2;3,3,1).
\end{gathered}
\]
\end{corollary}
\begin{proof}
The first and last satisfy the orbit lemma directly. The
second and third use the order-two endpoint corollary, and
the fourth uses the three-element sink lemma. In the frozen
original linear numbering these are $252,253,254,255,378$.
\end{proof}
